\section{How to Use This Template}


\subsection{Personalizing the template}
Edit file \verb|./setup/metadata.tex| to personalize the document. Here
you can enter your name and the title of the thesis, as well as the submission
date. Refer to \cref{sc:changingMainDocumentLanguage} for customizing the main
document language.

To select an appropriate logo, modify \verb|\titlePageLogo|. Following options
are currently available:
\begin{itemize}
  \item \verb|logo-none| : no logo;
  \item \verb|logo-mor| : MOR Lab logo;
  \item \verb|logo-mpc| : MPC Lab logo.
\end{itemize}


\subsection{Equations}
\subsubsection{Basics}
A simple equation:
\begin{equation}
  \label{eq:simple}
  2 + 2 = 4.
\end{equation}

Align and number separately:
\begin{align}
  \label{eq:align:separate1}
  2 + 2 &= 4; \\
  \label{eq:align:separate2}
  20 + 20 &= 40.
\end{align}

Align and number as subequations:
\begin{subequations}
  \label{eq:subequations}
  \begin{align}
    \label{eq:subequations:part1}
    2 + 2 &= 4; \\
    \label{eq:subequations:part2}
    20 + 20 &= 40.
  \end{align}
\end{subequations}

Referencing: \cref{eq:subequations} is the same as
\cref{eq:align:separate1} and \cref{eq:align:separate2},
but we can also reference the subequations:
\cref{eq:subequations:part1} and \cref{eq:subequations:part2}.

\subsubsection{Highlighting equations}
You can also highlight equations using the \verb|empheq| package.
Although the usage syntax is quite complex, you will get used to it
quickly:
\begin{empheq}[box = \bluebox]{equation}
  \pi \approx \frac{22}{7}.
\end{empheq}
By the way, notice that you can typeset bold Greek letters using the \verb|\T{}| macro:
\begin{empheq}[box = \orangebox]{align}
  \dot{\T{x}} &= \T{A} \T{x} + \T{B} \T{u}; \\[2 pt]
  \dot{\T{\xi}} &= \T{\Lambda} \T{\xi} + \T{B}_{\T{\xi}} \T{u}.
\end{empheq}
And, of course, \verb|empheq| supports subequations:
\begin{subequations}
  \begin{empheq}[box = \greenbox]{align}
    \dot{\T{x}} &= \T{A} \T{x} + \T{B} \T{u}; \\[2 pt]
    \T{y} &= \T{C} \T{x} + \T{D} \T{u}.
  \end{empheq}
\end{subequations}

\subsubsection{Units}
Use the \verb|siunitx| package for typesetting units:
\begin{subequations}
  \begin{align}
    c &= \SI{2.9979e8}{\metre\per\second}; \\
    \hbar &= \SI{6.63e-34}{\joule\second}; \\
    Z_\text{0} &= \SI{376.3}{\ohm}; \\
    \alpha^{-1} &\approx \num{137}.
  \end{align}
\end{subequations}
\verb|siunitx| also allows typesetting of numbers (\verb|\num{}|)
and units only (\verb|\si{}|). Trailing zeros are preserved; this
is very useful for typesetting tables (cf.\ \cref{tab:nums}).
\begin{table}
  \centering
  \caption[Example for nice number formatting]{%
    Example for nice number formatting (iteration %
    $x_{i+1} = x_i - \tan{x_i}$)}
  \label{tab:nums}
  \begin{tabular}{cr}
    \hline
    Iteration &      \TABC{Value} \\ \hline
      0     & \num{2.000000000} \\
      1     & \num{4.185039863} \\
      2     & \num{2.467893675} \\
      3     & \num{3.266186286} \\
      4     & \num{3.140943912} \\
      5     & \num{3.141592654} \\ \hline
  \end{tabular}
\end{table}


\subsection{Theorems and other statements}
\begin{dfn}[Optional definition description]
  This is a simple definition.
\end{dfn}

\begin{prp}[Optional proposition description]
  This is a simple proposition.
\end{prp}

\begin{lem}[Optional lemma description]
  This is a simple lemma.
\end{lem}

\begin{thm}[Optional theorem name]
  This is a simple theorem.
\end{thm}
\begin{proof}
  And this is its proof.
\end{proof}

\begin{cor}[Optional corollary description]
  This is a simple corollary.
\end{cor}

\begin{rem}[Optional remark description]
  This is a simple remark.
\end{rem}

\begin{exa}[Optional example description]
  This is a simple example.
\end{exa}


\subsection{Algorithms}
\cref{al:thesis} describes a method for writing high-quality theses.
You can reference a line in an algorithm by adding a \verb|\label{}|
after it, e.g.\ \cref{al:thesis:workMore} is essential.

\begin{algorithm}
  \setstretch{1.30}
  \KwData{task description, project duration $N$ (days)}
  \KwResult{thesis}
  \Begin{
    read literature\;
    \For{$i = 1\;:\;1\;:\;N$}{
      drink coffee\Comment*[r]{very important!}
      \While{problems unsolved}{
        solve problems\;
      }
      \eIf{tired}{
        go home\;
      }{
        work more\;\label{al:thesis:workMore}
      }
    }
    write paper\;
    submit paper\;
  }
  \caption{How to write a thesis}
  \label{al:thesis}
\end{algorithm}


\subsection{Figures}
You can reference the whole figure (\cref{fig:cuteAnimals})
as well as each subfigure separately:
\begin{itemize}
  \item \cref{fig:cuteAnimals:kitten1} was downloaded from \url{https://commons.wikimedia.org/wiki/File:Green_eyes_kitten.jpg}
  \item \cref{fig:cuteAnimals:kitten2} was downloaded from \url{https://commons.wikimedia.org/wiki/File:Golden_tabby_and_white_kitten_n03.jpg}
  \item \cref{fig:cuteAnimals:puppy1} was downloaded from \url{https://commons.wikimedia.org/wiki/File:Westie_pups.jpg}
  \item \cref{fig:cuteAnimals:puppy1} was downloaded from \url{https://commons.wikimedia.org/wiki/File:St._Bernard_puppy.jpg}
\end{itemize}

\begin{figure}[p]
  \centering
	\subfloat[green-eyed kitten]{\centering
	    \includegraphics[width = 0.45\linewidth]{./images/kitten1}
	    \label{fig:cuteAnimals:kitten1}}
	\subfloat[golden tabby and white kitten]{\centering
	    \includegraphics[width = 0.45\linewidth]{./images/kitten2}
	    \label{fig:cuteAnimals:kitten2}}\\[2em]
	\subfloat[Westie puppies]{\centering
	    \includegraphics[width = 0.45\linewidth]{./images/puppy1}
	    \label{fig:cuteAnimals:puppy1}}
	\subfloat[St.\ Bernard puppy]{\centering
	    \includegraphics[width = 0.45\linewidth]{./images/puppy2}
	    \label{fig:cuteAnimals:puppy2}}
   \caption[Cute animals]{This is an example for subfigures and cute animals}
  \label{fig:cuteAnimals}
\end{figure}

\subsection{Tikz}

\subsubsection{Drawing with Tikz}

You can also add and resize TikZ figures, e.g.\ \cref{fig:tikzExample}.

\begin{figure}[p]
	\centering
	\resizebox{0.4\linewidth}{!}{\begin{tikzpicture}
    \foreach \i in {0,1,...,15} \draw[thick] (0.0, \i/7.5) -- +(-0.1, 0.1);
    \draw[thick] (0, 0) -- +(0, 2);

    \draw[tum green, -latex'] (2.0, 3) -- +(1, 0) node[midway, above] {$x$};
    \draw[tum green] (3, 1.8) -- +(0, 1.3);
    \draw[tum green] (2.0, 3) +(0, -0.2) -- +(0, 0.2);
    \foreach \i in {0,1,...,3} \draw[tum green] (2.0, 2.8 + \i/7.5) -- +(-0.1, 0.1);

    \def\dampY{0.60}
    \def\dampLX{0}
    \def\dampRX{3}
    \def\dampLRod{1.25}
    \def\dampRRod{1.2}
    \def\dampWPiston{0.14}
    \def\dampWCylinder{0.20}
    \def\dampLCylinder{0.6}

    \draw[thick] (\dampLX, \dampY) -- +(\dampLRod, 0); % Left rod to the piston
    \draw[thick] (\dampLX + \dampLRod, \dampY + \dampWPiston/2) -- +(0, -\dampWPiston); % Piston
    \draw[thick] (\dampRX - \dampRRod - \dampWCylinder/2, \dampY - \dampWCylinder/2) arc (-90:+90:\dampWCylinder/2); % Cylinder head
    \draw[thick] (\dampRX - \dampRRod - \dampWCylinder/2, \dampY + \dampWCylinder/2) -- +(-\dampLCylinder, 0);  % Upper cylinder
    \node[above] at (\dampRX - \dampRRod - \dampWCylinder/2 - \dampLCylinder/2, \dampY + \dampWCylinder/2) {$d$};
    \draw[thick] (\dampRX - \dampRRod - \dampWCylinder/2, \dampY - \dampWCylinder/2) -- +(-\dampLCylinder, 0); % Lower cylinder
    \draw[thick] (\dampRX, \dampY) -- +(-\dampRRod, 0); % Right rod to the cylinder head

    \def\springY{1.4}
    \def\springLX{0}
    \def\springRX{3}
    \def\springLRod{0.5}
    \def\springRRod{0.5}
    \def\springWCoil{0.20}
    \def\springNCoil{12}
    \pgfmathsetmacro{\springLCoil}{(\springRX - \springRRod - \springLX - \springLRod)/\springNCoil}

    \draw[thick] (\springLX, \springY) -- ++(\springLRod + 0.01, 0); % Left rod
    \draw[thick] (\springRX, \springY) -- ++(-\springRRod - 0.01, 0); % Right rod
    \foreach \i in {1, 2, ..., \springNCoil} %
        \draw[thick] (\springLX + \springLRod + \i*\springLCoil - \springLCoil, \springY) -- %
        ++(\springLCoil/4, -\springWCoil/2) -- ++(\springLCoil/2, +\springWCoil) -- ++(\springLCoil/4, -\springWCoil/2);
    \node[above] at ( \springLX/2 + \springLRod/2 + \springRX/2 - \springRRod/2, \springY + \springWCoil/2) {$c$};

    \draw[thick] (3, 0.2) rectangle +(2, 1.6) node[midway] {$m$};
\end{tikzpicture}
}
	\caption{Example for a TikZ figure}
	\label{fig:tikzExample}
\end{figure}

\begin{figure}[ht]
	\centering
	\begin{tikzpicture}

% Stützstellen
\draw [tum blue 3, fill]  (-3, -0.05) rectangle (-1, 0.05); 

\draw (-3, 0) node[below=0.6em] {$\underline{\mathbf{p}}_{i-1}$}  +(-0.2, 0.2)  -- +(0.2, -0.2);
\draw (-3, 0) node[above=0.6em] {$\underline{\mathbf{V}}_{i-1}$}  +(-0.2, -0.2)  -- +(0.2, 0.2);

\draw (-1, 0) node[below=0.6em] {$\overline{\mathbf{p}}_{i-1}$}  +(-0.2, 0.2)  -- +(0.2, -0.2);
\draw (-1, 0) node[above=0.6em] {$\overline{\mathbf{V}}_{i-1}$}  +(-0.2, -0.2)  -- +(0.2, 0.2);

%-------------------------------------------------------------------------------------------------
\draw [tum blue 2, fill]  (1, -0.05) rectangle (3, 0.05); 

\draw (1, 0) node[below=0.6em] {$\underline{\mathbf{p}}_{i}$}  +(-0.2, 0.2)  -- +(0.2, -0.2);
\draw (1, 0) node[above=0.6em] {$\underline{\mathbf{V}}_{i}$}  +(-0.2, -0.2)  -- +(0.2, 0.2);

\draw (3, 0) node[below=0.6em] {$\overline{\mathbf{p}}_{i}$}   +(-0.2, 0.2)  -- +(0.2, -0.2);
\draw (3, 0) node[above=0.6em] {$\overline{\mathbf{V}}_{i}$}  +(-0.2, -0.2)  -- +(0.2, 0.2);

% x-Achse
\draw[-latex'] (-4, 0) -- (4, 0) node[below] {$\mathbf{p}$};

% V_{t-1}
\draw[tum blue 4] (-2, 0) node[below=0.1em] {$\mathbf{p}_{t-1}$};
\draw[tum blue 4, thick, latex'-] (-2,0) +(0, 0) -- +(0, 2) node[above] {$\mathbf{V}_{t-1}$}; 

% V_{t}
\draw[tum blue 3] (2, 0) node[below=0.1em] {$\mathbf{p}_t$}; 
\draw[tum blue 3, thick, latex'-] (2,0) +(0, 0) -- +(0, 2) node[above] {$\mathbf{V}_{t}$}; 

% Dashed line
\draw[dashed,->] (-2, 1.5) -- (2, 1.5);

\end{tikzpicture}
	\caption{Illustration der verschiedenen Parameterstuetzstellen und deren lokalen Basen}
	\label{fig:ableitung_V}
\end{figure} 

\begin{figure}[hbp]
	\centering
	\begin{tikzpicture}
	 \node[anchor=south west,inner sep=0] (image) at (0,0) {\includegraphics[width = 10 cm]{./images/wagen}};
	    \begin{scope}[x={(image.south east)},y={(image.north west)}]
	        \draw[white, thick, latex'-] (0.63, 0.74) -- +(0.1, 0.1) node[right] {Wagen};
	        \draw[white, thick, latex'-] (0.25, 0.58) -- +(-0.05, 0.1) node[left] {Schiene};
	        \draw[white, thick, latex'-] (0.32, 0.42) -- +(-0.05, -0.15) node[left, fill=black] {Sensorrad};
	        \draw[white, thick, latex'-] (0.49, 0.37) -- +(0.05, -0.15) node[right, fill=black] {Antriebsrad};
	    \end{scope}
\end{tikzpicture}
	\caption{Wagen in der Schienenfuehrung}
	\label{fig:aufbau:wagen}
\end{figure}

\begin{figure}
	\centering
	\begin{minipage}[c]{\textwidth}
		\centering
		%cwt-a
\begin{tikzpicture}[domain=0:3, scale=0.48]
\centering
		%KoordinatenSystem
		\draw[->] (0,0) -- (17,0) node[right]{$t$};
		\draw[->] (0,0) -- (0,17) node[left]{$a$};
		
		%-Rand vom Kasten
		\draw[] (16,0) -- (16,16);
		\draw[] (0,16) -- (16,16);
		
		%waagrechte Striche
		\draw[] (0,15) -- (16,15);
		\draw[] (0,14) -- (16,14);
		
		\draw[] (0,12) -- (16,12);

		\draw[] (0,8) -- (16,8);
		
		%senkrechte Striche		
		\draw[] (2,0) -- (2,8);
		\draw[] (4,0) -- (4,12);
		\draw[] (6,0) -- (6,8);
		\draw[] (8,0) -- (8,14);
		\draw[] (10,0) -- (10,8);
		\draw[] (12,0) -- (12,12);
		\draw[] (14,0) -- (14,8);			
\end{tikzpicture}




	
		\caption{Zeit-Skalierungsebene bei der Wavelet-Transformation}
		\label{fig:cwt_a}
	\end{minipage}\\[1em]
	\begin{minipage}[c]{\textwidth}
		\centering
		%cwt-f
\begin{tikzpicture}[domain=0:3, scale=0.48]
\centering
		%KoordinatenSystem
		\draw[->] (0,0) -- (17,0) node[right]{$t$};
		\draw[->] (0,0) -- (0,17) node[left]{$f$};
		
		%-Rand vom Kasten
		\draw[] (16,0) -- (16,16);
		\draw[] (0,16) -- (16,16);
		
		%waagrechte Striche
		\draw[] (0,1) -- (16,1);
		\draw[] (0,2) -- (16,2);
		\draw[] (0,4) -- (16,4);
		\draw[] (0,8) -- (16,8);
		\draw[] (0,16) -- (16,16);
		
		%senkrechte Striche		
		\draw[] (2,16) -- (2,8);
		\draw[] (4,16) -- (4,4);
		\draw[] (6,16) -- (6,8);
		\draw[] (8,16) -- (8,2);
		\draw[] (10,16) -- (10,8);
		\draw[] (12,16) -- (12,4);
		\draw[] (14,16) -- (14,8);
		
\end{tikzpicture}
		\caption{Zeit-Frequenzebene bei der Wavelet-Transformation}
		\label{fig:cwt_f}
	\end{minipage}
\end{figure}

%===================================================
\subsubsection{Matlab2Tikz}
You can download matlab2tikz from \url{https://de.mathworks.com/matlabcentral/fileexchange/22022-matlab2tikz-matlab2tikz}.

\begin{figure}[t!]
	\centering
	\begin{minipage}{0.95\linewidth}
		\centering
		\setlength\figureheight{3.3cm}
		\setlength\figurewidth{0.88\textwidth}
		\input{images/matlab/prbs.tikz}
	\end{minipage}
	\\[0.6em]
	\begin{minipage}{0.95\linewidth}
		\centering
		\setlength\figureheight{3.3cm}
		\setlength\figurewidth{0.88\textwidth}
		\input{images/matlab/prbs_LDS_nicht_dB.tikz}
	\end{minipage}
	\\[0.3em]
	\caption[PRBS-Anregungssignal mit zugehoerigem Leistungsdichtespektrum]{Darstellung des PRBS-Anregungssignals mit zugehoerigem Leistungsdichtespektrum.}
	\label{fig:prbs}
\end{figure}

\begin{figure}[tp!]
	\centering
	\setlength\figureheight{4cm}
	\setlength\figurewidth{0.8\textwidth}
	\input{images/matlab/aufloesung_fft.tikz}
	\caption[Aufloesung im Zeitbereich und im Frequenzbereich]{Aufloesung im Zeitbereich (oben) und im Frequenzbereich (unten)}
	\label{fig:aufloesung_fft}
\end{figure}

\begin{figure}[hbp!]
	\centering
	\subfloat[$v = 1 \frac{m}{s}$]{\centering
		\setlength\figureheight{5cm}
		\setlength\figurewidth{6.3cm}
		\input{images/matlab/L_1_v_1_r_10.tikz}		
		\label{fig:L_1_v_1_r_10}}
	\subfloat[$v = 5 \frac{m}{s}$]{\centering
		\setlength\figureheight{5cm}
		\setlength\figurewidth{6.3cm}
		\input{images/matlab/L_1_v_5_r_10.tikz}		
		\label{fig:L_1_v_5_r_10}}\\
	\subfloat[$v = 10 \frac{m}{s}$]{\centering
		\setlength\figureheight{5cm}
		\setlength\figurewidth{6.3cm}
		\input{images/matlab/L_1_v_10_r_10.tikz}		
		\label{fig:L_1_v_10_r_10}}
	\subfloat[$v = 20 \frac{m}{s}$]{\centering
		\setlength\figureheight{5cm}
		\setlength\figurewidth{6.3cm}
		\input{images/matlab/L_1_v_20_r_10.tikz}		
		\label{fig:L_1_v_20_r_10}}
	\caption{Simulationsergebnisse bei einseitiger Reduktion mit Eingangs-Krylow-Raum}
	\label{fig:red_r_10_input_Krylov}
\end{figure}

%===================================================
\subsubsection{Inkscape}

\subsubsection{Inkscape documents saved as .eps}

\begin{figure}[htbp]
	\centering
	\includegraphics[width=0.7\textwidth]{images/inkscape/circuits/12V_Bordnetz}
	\caption{Prinzipschaltbild des bereits etablierten 12V-Bordnetzes}
	\label{fig:12V_Bordnetz}
\end{figure}

\begin{figure}[htbp]
	\centering
	\includegraphics[width=1.0\textwidth]{images/inkscape/circuits/48V_12V_Bordnetz}
	\caption{Prinzipschaltbild des zukünftigen 48V/12V-Bordnetzes mit UltraCap-Modul und DC/DC-Wandler}
	\label{fig:48V_12V_Bordnetz}
\end{figure}

\subsubsection{Inkscape document with pdf-tex}
Just save the inkscape document as pdf and choose the option \emph{PDF+LaTeX: Text in PDF weglassen und LaTeX Datei erstellen}.

\begin{figure}[ht]
	\begin{center}
		\def\svgwidth{\linewidth}\small
		\import{./images/inkscape/charts/}{flowchart.pdf_tex}
		\caption{Overview of work packages and assignments}
		\label{fig:assignment_wp}
	\end{center}
\end{figure}

\begin{figure}[ht]
	\centering
	\def\svgwidth{\linewidth}\footnotesize
	\import{./images/inkscape/charts/}{gantt_chart.pdf_tex}
	\caption{Gantt chart of the proposed work program}
	\label{fig:Gantt}
\end{figure}

%===================================================
\subsection{Blockschaltbilder}
You can draw block diagrams with Tikz. For this, you can use the commands defined in the TeX-File \emph{blockschaltbilder} (setup/blockschaltbilder). Some examples follow now: see Fig.~\ref{fig:zwei-freiheitsgrade-regelung}, Fig,~\ref{fig:zwei-freiheitsgrade-regelung-zustand}.

\begin{figure}[tp]
	\centering
	\scalebox{0.95}{
		\input{tikz/zwei-freiheitsgrade-regelung}}
	\caption{Something interesting.}
	\label{fig:zwei-freiheitsgrade-regelung}
\end{figure}

\begin{figure}[tp]
	\centering
	\scalebox{0.95}{
		\input{tikz/zwei-freiheitsgrade-regelung-zustand}}
	\caption{Something interesting.}
	\label{fig:zwei-freiheitsgrade-regelung-zustand}
\end{figure}

	

%===================================================
\subsection{ToDo Notes}
%\tikzexternaldisable
\todoX{Read the newspaper and the magazine}

\todoY{Make the breakfast}

\listoftodos
%\tikzexternalenable


%===================================================
\subsection{Tables}
\LaTeX{} tables are trivial! Cf.\ \cref{tab:example}.

\begin{table}
  \centering
  \caption{Funny webcomics}
  \label{tab:example}
  \begin{tabular}{llll}
    \hline
    Title                    & Author     & Date         & Link                                                    \\
    \hline
    Highway Engineer Pranks  & R.\ Munroe & 2007-04-25   & \href{http://xkcd.com/253/}{xkcd.com}                   \\
    Stop Being Engineers!    & S.\ Adams  & 2014-02-02   & \href{http://dilbert.com/strip/2014-02-02}{dilbert.com} \\
    The Typographical Terror & D.\ Malki  & 2010--08--20 & \href{http://wondermark.com/650/}{wondermark.com}       \\
    \hline
  \end{tabular}
\end{table}


\subsection{Listings}
Quine in \cref{lst:quine:c} was found \href{https://www.nyx.net/~gthompso/self_c.txt}{here} (search for Jeff Hollingsworth's quine),
quine in \cref{lst:quine:matlab} \href{http://blog.garritys.org/2009/10/quines.html}{here}. The code in
\cref{lst:sierpinski} was found \href{http://rosettacode.org/wiki/Sierpinski_carpet#Python}{here}.

\begin{lstlisting}[language = C, label = lst:quine:c, caption = {An elegant C quine}]
#define Z(q)q,#q
main()printf(Z("#define Z(q)q,#q\nmain()printf(Z(%s));\n"));
\end{lstlisting}

\begin{lstlisting}[label = lst:quine:matlab, caption = {An elegant MATLAB quine}]
a=['ns}z2e1kGe1116k6111gE16;:6ek7;:61gg3E1'];
disp(['a=[''',a,'''];',10,[a-10,']]);']]);
\end{lstlisting}

\lstinputlisting[language = Python, label = lst:sierpinski, caption = {Sierpinski carpet in Python}]{./sources/sierpinski.py}


\subsection{Languages support}
\subsubsection{Typesetting non-English texts}
Use the \verb|otherlanguage| environment to typeset non-English texts:
\begin{center}
  \fbox{%
    \begin{minipage}{0.7\linewidth}
      \begin{otherlanguage}{ngerman}
        Franz jagt im komplett verwahrlosten
        Taxi quer durch Bayern. Zwölf Boxkämpfer jagen Viktor
        quer über den großen Sylter Deich. Vogel Quax zwickt
        Johnys Pferd Bim. Sylvia wagt quick den Jux bei
        Pforzheim. Polyfon zwitschernd aßen Mäxchens Vögel
        Rüben, Joghurt und Quark. \glqq Fix, Schwyz!\grqq\ quäkt
        Jürgen blöd vom Paß. Victor jagt zwölf Boxkämpfer quer
        über den großen Sylter Deich. Falsches Üben von
        Xylophonmusik quält jeden größeren Zwerg.
        Heizölrückstoßabdämpfung.
      \end{otherlanguage}
    \end{minipage}
  }
\end{center}

\subsubsection{Changing the main document language}
\label{sc:changingMainDocumentLanguage}
To use German as a default language, edit the first line in \verb|_thesis.tex|:
Change\\
\verb|\documentclass[<...>, ngerman, english]{book}|\\
to\\
\verb|\documentclass[<...>, english, ngerman]{book}|.

You can also use \verb|british| instead of \verb|english| if you prefer BrE.


\subsection{Abbreviations and acronyms}
Abbreviations and acronyms are defined in \verb|auxfiles/abbreviationsAndAcronyms.tex|.
You can typeset them using the \verb|\ac{}| command. Refer to \texttt{acro} package
documentation for specialized commands (\verb|\acs{}|, \verb|\acl{}|, \verb|\acp{}|,
\verb|\iac{}|, \verb|\Iac{}|, \ldots).

Usage example:
\begin{itemize}
  \item An abbreviation is explained if used for the first time: \ac{MPC};
  \item subsequently, only the short version is used: \ac{MPC};
  \item however, you can force \texttt{acro} to use the long version anytime: \acl{MPC}.
\end{itemize}

\emph{Important:} If you want to have only the actually used abbreviations and acronyms
in the list (\cref{sc:abbreviationsAndAcronyms}), edit \verb|setup/settings.tex|: Find
the \verb|\acsetup{}| command and set \verb|list/display| to \texttt{used}.


\subsection{References}
\subsubsection{Hyperlinks}
See \url{https://en.wikibooks.org/wiki/LaTeX/Hyperlinks}.

Also, always use \verb|\cref{}| for referencing equations, tables or other things!

\subsubsection{Literature}
Cite multiple references:
\textcite{Bronstein2012,Follinger2013,Gamma2004,Gantmacher2000a,Gantmacher2000b,%
  Golub2013,Horn2013,LaSalle1961,Mehlhorn2008,Sontag1998,Adamy2009,Horn1994,Watkins2002}.

If you want to have only the reference number:
\cite{Bronstein2012,Follinger2013,Gamma2004,Gantmacher2000a,Gantmacher2000b,Golub2013,%
  Horn2013,LaSalle1961,Mehlhorn2008,Sontag1998,Adamy2009,Horn1994,Watkins2002}.

If you want to reference a specific section: \cite[\S~5.2.2, pp.~248--250]{Golub2013}.

If you want to reference a specific equation: \cite[Eq.~(13.22), p.~721]{Bronstein2012}.
